\documentclass[11pt]{article}

\usepackage[margin=1.1in]{geometry}
\usepackage{times}
\usepackage[T1]{fontenc}
\usepackage[utf8]{inputenc}
\usepackage{microtype}
\usepackage{amsmath}
\usepackage{amssymb}
\usepackage{booktabs}
\usepackage[round]{natbib}
\usepackage{url}

\newcommand{\bl}{b_L}
\newcommand{\blh}{\hat{b}_L}
\newcommand{\wg}{w_{\mathrm{g}}}
\newcommand{\wv}{w_{\mathrm{v}}}


\title{The Carried Reference Tunes a Switch--Multiplier\\Continuum in Gated (GLU) MLPs\\
\large A companion note to \emph{Transformer MLP Gate Thresholds Are Couplings
to a Carried Reference Direction}}
\author{Olli Tuomi\\Evident Solutions Oy}
\date{September 2026}

\begin{document}
\maketitle

\begin{abstract}
A prior paper showed that a trained transformer implements its MLP gate
thresholds as a coupling $w\cdot\bl$ to a carried reference direction $\bl$, the
corpus-mean layer-normalised gate input. We use that reference to give a
functional reading of the gate in a gated (GLU) MLP. Each gate unit's coupling to
$\bl$ places it on a continuous \emph{switch--multiplier} spectrum: a unit
strongly (negatively) coupled to $\bl$ rests closed and behaves as a thresholded
switch, while a unit orthogonal to $\bl$ rests at the activation's knee and, to
first order, computes a bilinear product of its two branch reads. The per-unit
cosine to $\bl$ is the coordinate along that spectrum, and trained SwiGLU
allocates units continuously across it, shifting toward the multiplier end with
depth. Supporting this, the coupling localises to the gate branch (across four
SwiGLU models the \texttt{gate\_proj} rows align with $\bl$ at $4$--$33\times$ the
\texttt{up\_proj} rows) and lives in the reference's high-variance subspace; a
bit-identical-init training ladder shows the elementwise nonlinearity is what
creates the switch end, and that among four feed-forward variants only the gated
one spans both ends with independent key and value reads. We suggest this
coverage, rather than any single mode, as a reason the GLU family is an effective
default, and note the reading extends to the experts of a mixture-of-experts. The
account is geometric and makes no quality claim.
\end{abstract}

\section{Introduction}
\label{sec:intro}

A companion paper \citep{tuomi2026background} reported that a trained
transformer implements its MLP gate thresholds as a coupling to a carried
reference direction. The reference $\bl$ is the corpus mean of the
layer-normalised MLP input, the vector the gate rows read; a gate's resting
pre-activation decomposes exactly as $\langle w\cdot x_{\mathrm{ln}}\rangle + b_n
= w\cdot\bl + b_n$, and in the models studied the coupling $w\cdot\bl$ carries
the resting inhibition of more than $99.9\%$ of gates, at $50$--$60\times$ the
explicit bias parameter. Per read row the coupling is a thin slice of the row's
norm: factoring each row as $w = c_n\,\blh + w_\perp$ (the anchor paper's own
notation), the fraction $|c_n|/\lVert w\rVert$ has median $0.07$--$0.12$, so about $99\%$ of every row is
left free for content. We take this decomposition as established and build on it.

That paper reads the coupling as a threshold. A gated MLP invites a functional
question the single-matrix picture does not: its read stage factorises into a
gate branch (carrying the elementwise nonlinearity) and a linear value branch,
whose product is written out, so a unit's coupling to $\bl$ can vary within a
layer. We find it varies along one axis. Expanding the SiLU at its knee, a gate
unit whose operating point $\wg\!\cdot\!\bl$ is deep-negative rests closed and
fires only above threshold, a \emph{switch}; a unit whose operating point sits at
the knee ($\wg \perp \bl$) computes, to leading order,
$\tfrac{1}{2}(\wg\!\cdot\! x)(\wv\!\cdot\! x)$, a bilinear \emph{multiplier} of its
two reads. A unit's cosine to $\bl$ is its coordinate on this switch--multiplier
spectrum; trained SwiGLU spreads units continuously along it, and the mix moves
toward the multiplier end with depth. The rest of the note supports and bounds
this reading: the coupling localises to the gate branch and to the reference's
high-variance subspace (\S\ref{sec:localize}); a controlled ladder shows the
nonlinearity is what creates the switch end, and that of four feed-forward
variants only the gated one spans the spectrum with independent key and value
reads (\S\ref{sec:ladder}); freeing the reference from the value branch returns a
content dimension (\S\ref{sec:capacity}); and we discuss why this coverage suits
the GLU family, and how the reading should extend inside a mixture-of-experts
(\S\ref{sec:discussion}).

The claim is geometric, about where a
shared rank-one direction sits in the read weights, and not a claim about
quality: whether gating improves the loss is a separate question, addressed
elsewhere at this feed-forward geometry \citep{shazeer2020}, and we make no
quality claim here. And what localises to the gate is the shared reference slice
only: both branches carry content, so
that freezing either one at its resting value does comparable damage to prediction
(\S\ref{sec:scope}), as the product of two factors requires.

Throughout, ``gate'' is the branch carrying the elementwise nonlinearity
($\wg$, \texttt{gate\_proj}), ``value'' the linear branch ($\wv$,
\texttt{up\_proj}); a SwiGLU MLP writes $W_{\mathrm{down}}\!\left(\mathrm{SiLU}
(\wg x_{\mathrm{ln}})\odot(\wv x_{\mathrm{ln}})\right)$. We use the reference
$\bl := \langle x_{\mathrm{ln}}\rangle$ of the anchor paper: the corpus-mean
layer-normalised MLP input, the direction the gates read. In the RMSNorm models
$x_{\mathrm{ln}}$ is the RMS-normalised input (no mean subtraction); $\bl$ is
well defined and is still what the gate rows dot.

\section{The reference coupling places a unit on a switch--multiplier spectrum}
\label{sec:spectrum}

A SwiGLU unit writes $\mathrm{SiLU}(g)\,v\,d$ with $g=\wg\!\cdot\! x$,
$v=\wv\!\cdot\! x$, and $d$ its down-projection column. Expanding the SiLU about
its knee, $\mathrm{SiLU}(z)\approx \tfrac12 z + \tfrac14 z^2$, gives two regimes
selected by where the unit rests, $g_0=\wg\!\cdot\!\bl$:
\begin{itemize}
\item \textbf{Switch.} $g_0\ll 0$: the gate rests closed, $\mathrm{SiLU}(g)\approx
  0$, and the unit fires only when the input pushes $g$ above threshold, writing
  its value when it does. A thresholded, off-by-default detector.
\item \textbf{Multiplier.} $g_0\approx 0$ ($\wg \perp \bl$): the gate rests at the
  knee, $\mathrm{SiLU}(g)\,v \approx \tfrac12(\wg\!\cdot\! x)(\wv\!\cdot\! x)$, an
  always-on bilinear product of the two branch reads. This is the object a
  bilinear MLP is built from \citep{pearce2024}.
\end{itemize}
The operating point $g_0$, equivalently a unit's cosine to $\bl$, is therefore a
coordinate that slides the unit between the two regimes. (This is a per-unit
coordinate set by the carried reference, not the global activation shape:
SiLU/Swish interpolate linear-to-ReLU through a shared $\beta$
\citep{ramachandran2017swish}, whereas here each unit's regime is set by its own
alignment to $\bl$.)

Trained SwiGLU populates the whole spectrum. Table~\ref{tab:modes} splits gate
units at three layers of Qwen2.5-3B into the switch end (coupling above
$2\times$ the random-alignment floor) and the multiplier end (below it), and
reports each group's \emph{duty} -- the anchor paper's term for the fraction of
positions at which a unit's pre-activation is positive. The switch units rest
deep and fire rarely (duty near $0$), with a pre-activation dominated by the
constant operating point (ratio of $|{\rm mean}|$ to standard deviation well above
one); the multiplier units rest near the knee, fire about a third of the time
(duty $\approx 1/3$), and are input-driven (that ratio near zero). The coupling is graded rather than
bimodal (the cosine distribution is a smooth tail, no gap), predominantly
inhibitory in sign ($88$--$100\%$ of gate units), and scales with the row norm
(Spearman $0.55$--$0.66$); the multiplier fraction rises with depth ($4\%\to
31\%\to 39\%$ here). The multiplier-end units are input-driven content gates, not
a second reference: their rows share no dominant common axis (top singular value
$\leq 1.6\%$ of energy) and are not aligned with the value branch
($\cos(\wg,\wv)\approx 0$).

These labels are behavioural, not only geometric. Comparing each unit's actual
output $\mathrm{SiLU}(g)\,v$ to the leading bilinear term $g\,v$ on tokens, the
knee (uncoupled) units are well approximated by the product -- median $R^2$
$0.55$--$0.83$ and best-fit slope $\approx 0.4$, near the knee's coefficient of
$\tfrac12$ -- so their output genuinely is the bilinear term. The deeply-coupled
switch units are not: at the shallowest layer, where switches rest deepest, their
$R^2$ falls to $0.15$ with slope $\approx 0$, the behaviour of a threshold rather
than a product. The bilinear fidelity decreases with coupling at every layer
(Spearman between $R^2$ and $|\cos(\wg,\blh)|$ is $-0.17$ to $-0.42$), so a unit's
geometric coordinate predicts whether it actually computes a product. The gap is
sharp at the extremes and graded between them, where softly-resting units track
the product either way -- the signature of a continuum rather than two classes.

\begin{table}[t]
\centering\small
\begin{tabular}{llrrrr}
\toprule
$L$ & mode & share & rest $\wg\!\cdot\!\bl$ & duty & $|{\rm mean}|/{\rm sd}$ \\
\midrule
9  & switch     & 96\% & $-4.79$ & 0.00 & 3.2 \\
   & multiplier &  4\% & $-0.73$ & 0.19 & 0.8 \\
18 & switch     & 69\% & $-0.62$ & 0.13 & 1.1 \\
   & multiplier & 31\% & $-0.22$ & 0.32 & 0.4 \\
27 & switch     & 61\% & $-0.87$ & 0.11 & 1.3 \\
   & multiplier & 39\% & $-0.23$ & 0.34 & 0.4 \\
\bottomrule
\end{tabular}
\caption{Gate units at the two ends of the spectrum (Qwen2.5-3B), split by
coupling to $\bl$. Switch units rest deep-negative, fire rarely, and are
constant-dominated; multiplier units rest near the knee, fire $\sim\!1/3$ of the
time, and are input-driven. The multiplier share grows with depth. Source:
\texttt{census\_gated\_coupling\_dist.py}, \texttt{census\_gated\_uncoupled\_units.py}.}
\label{tab:modes}
\end{table}

\section{The coupling localises to the gate branch}
\label{sec:localize}

The switch end of the spectrum lives in the gate branch, and its geometry is
visible in the weights alone. For each hidden unit we take the
cosine of its gate row $\wg$ and its value row $\wv$ with the unit reference
$\blh$, normalised by the row norm, and read the median over units at three
mid-stack layers of Qwen2.5-3B (Table~\ref{tab:geom}). The gate rows carry the
coupling and the value rows do not: $|\cos(\wg,\blh)|$ is $34\times$ the value
figure at L9 and $7$--$8\times$ deeper, a dissociation of the same quantity the
companion paper measured inside a single matrix, now split across the two
projections. The sign is consistent with a threshold: the gate's resting
pre-activation $\wg\!\cdot\!\bl$ is negative, so the resting gate is closed
($\mathrm{SiLU}(\wg\!\cdot\!\bl) < 0$ at every layer) and opens conditionally.
The two branches are mutually orthogonal ($\cos(\wg,\wv)\approx 0$), each reading
a distinct direction, and the gate rows place more of their energy on the
high-variance coordinates of $x_{\mathrm{ln}}$ than the value rows do
($0.03$--$0.04$ against $0.015$, against an isotropic $0.031$), the first sign
that the coupling lives in the amplitude part of the reference, which the control
below makes precise. A forward pass agrees, weakly: the gate activation tracks
$\lVert x_{\mathrm{ln}}\rVert$ across positions at $1.2$--$1.9\times$ the value's
correlation, strongest at L9.

\begin{table}[t]
\centering\small
\begin{tabular}{lrrrrr}
\toprule
$L$ & $|\cos(\wg,\blh)|$ & $|\cos(\wv,\blh)|$ & ratio &
  resting $\mathrm{SiLU}(\wg\!\cdot\!\bl)$ & $\cos(\wg,\wv)$ \\
\midrule
9  & 0.0799 & 0.00232 & 34.4 & $-0.042$ & 0.000 \\
18 & 0.047 & 0.0056 &  8.4 & $-0.180$ & 0.001 \\
27 & 0.044 & 0.0062 &  7.1 & $-0.204$ & $-0.001$ \\
\bottomrule
\end{tabular}
\caption{Gate vs.\ value read geometry, Qwen2.5-3B (weights and $\bl$ only,
mid-stack layers of 36). Gate rows align with the reference $7$--$34\times$ the
value rows; the gate rests closed; the two branches are mutually orthogonal
(not self-gating). Gate high-variance-coordinate energy exceeds the value's
throughout ($0.03$--$0.04$ vs.\ $0.015$ against an isotropic $0.031$).
Source: \texttt{census\_gated\_gate\_structure.py}, Leg B.}
\label{tab:geom}
\end{table}

\paragraph{Where in the reference the coupling sits.}
On Phi-2 the companion paper finds mid-stack $\bl$ to be a distributed direction,
distinct from the massive-activation and outlier-coordinate objects
\citep{tuomi2026background}. The RMSNorm SwiGLU models here are the other case:
$\bl$ concentrates, with the eight highest-variance coordinates carrying
$66$--$75\%$ of $\lVert\bl\rVert^2$ across the three layers. This concentration
reflects the model's massive-activation profile rather than its normalisation
scheme: across seven models of both normalisation types, $\bl$'s largest
coordinates coincide with the massive-activation coordinates and its
concentration tracks their strength, which is why the same reference reads as
distributed in Phi-2 and concentrated here \citep{sun2024}. Either way, whether
the gate reads the reference or merely those coordinates has to be settled again
for these models, and for the branch split in particular. Zeroing the eight
coordinates and renormalising, so that two-thirds of the reference is discarded,
the gate still aligns with what remains at $12.8 / 4.0 / 3.2\times$ the value
(L9/L18/L27) and still rests closed. The localisation narrows as more of the
high-variance subspace is removed: with the top sixty-four coordinates gone the
ratio falls to about $1.3$ and the value's alignment rises to meet the gate's in
the low-variance bulk. The gate carries the reference coupling within the
reference's high-variance subspace, the same subspace the energy reading above
points to.

\paragraph{Replication across SwiGLU lineages.}
The same static reading on three mid-stack layers of four SwiGLU models
(Table~\ref{tab:xmodel}) puts the coupling in the gate in every one. The gate
rows align with $\bl$ more than the value rows at all twelve model-layers, and
the gate rests closed (negative resting SiLU) everywhere; the strength varies by
lineage, with Qwen sharpest. The clearest soft spot is SmolLM2's deepest sampled
layer, where the ratio falls to $1.8$ and the gate opens at a third of positions
rather than a tenth. Splitting SmolLM2's own gate units by the same
switch/multiplier threshold as Table~\ref{tab:modes} confirms this is a
population shift, not an unexplained dip: the multiplier-end share rises
$30\%\to41\%\to60\%$ from the shallowest to the deepest sampled layer, the same
monotonic climb as Qwen2.5-3B's $4\%\to31\%\to39\%$ (\S\ref{sec:spectrum}), only
starting and ending further along the spectrum. Across the other eleven layers
the separation is unambiguous, so the localisation is a property of the SwiGLU
construction and not of one model, with each model's multiplier-end mixture
setting how far it softens.

\begin{table}[t]
\centering\small
\begin{tabular}{lrrrl}
\toprule
model & \multicolumn{3}{c}{gate/value $|\cos(\cdot,\blh)|$ ratio} & resting SiLU \\
      & $L{\approx}.25$ & $L{\approx}.5$ & $L{\approx}.75$ & \\
\midrule
Qwen2.5-1.5B   & 11.6 & 7.3 & 8.0 & all $<0$ \\
Qwen2.5-3B     & 33.0 & 8.1 & 7.1 & all $<0$ \\
TinyLlama-1.1B &  4.4 & 5.2 & 4.7 & all $<0$ \\
SmolLM2-1.7B   &  5.3 & 4.8 & \textbf{1.8} & all $<0$ \\
\bottomrule
\end{tabular}
\caption{Cross-model replication (weights $+\bl$, three mid-stack layers each).
The reference coupling sits in the gate branch in every model; the gate rests
closed everywhere. It weakens at SmolLM2's deepest sampled layer (ratio $1.8$,
gate frac-open $0.34$), the clearest soft spot; splitting SmolLM2's own gate
units by the switch/multiplier threshold of Table~\ref{tab:modes} confirms a
population shift toward the multiplier end there ($30\%\to41\%\to60\%$ from
shallowest to deepest sampled layer).
Sources: \texttt{census\_gated\_branch\_xmodel.py},
\texttt{census\_gated\_uncoupled\_smollm2.py}.}
\label{tab:xmodel}
\end{table}

\section{The nonlinearity creates the switch end, and only the gated arm spans the spectrum}
\label{sec:ladder}

The switch end of the spectrum is the part that needs a carried reference: a bare
product has no off-by-default point. To see that the elementwise nonlinearity is
what supplies it, and which architectures can reach both ends, we use a controlled
training ladder. Four feed-forward variants were trained at the 160M
(Pythia) scale on identical data in identical order under a shared schedule,
differing only in the feed-forward block: a plain GELU MLP, a SwiGLU, a bilinear
GLU (the product of two linear branches, with no elementwise nonlinearity), and a
squared-ReLU MLP. The matched-shape arms share bit-identical initialisation
(maximum absolute weight difference $0$): the SwiGLU gate and value rows are the
bilinear arm's two branches at initialisation, and the GELU read matrix is the
squared-ReLU arm's. The SwiGLU-versus-bilinear pair therefore isolates a single
edit, adding the SiLU to one branch, and the GELU-versus-squared-ReLU pair
isolates swapping the elementwise function, each with everything else fixed.

Taking the cross-arm cosine of $\bl$ at matched layers (Table~\ref{tab:reloc}),
adding the gate relocates the reference while swapping the elementwise function
leaves it in place: the SwiGLU--bilinear cosine falls to $0.17$ over training
while the GELU--squared-ReLU cosine holds at $0.74$.\footnote{$\bl$ is the mean
post-LayerNorm MLP input, matching the anchor's definition. Measuring the mean
pre-normalisation residual instead gives the same dissociation at a wider spread
($0.12$ versus $0.93$), so it is not an artifact of the reference's definition.}
Both pairs begin aligned at initialisation ($0.89$ and $0.97$), so the divergence
is trained rather than a starting condition. The mechanism is the one the mature
models show: over the same training the SwiGLU gate develops its gate-over-value
alignment ratio (rising to about $13$ mid-stack) and its resting gate closes,
while the bilinear arm's two symmetric branches stay at ratio $1$ with no closed
rest, a built-in null that differs only by the SiLU. The SiLU places the coupling
in the gate, and placing it there is what moves the reference.

\begin{table}[t]
\centering\small
\begin{tabular}{lrrr}
\toprule
pair (edit) & step 0 & step 12k & step 150k \\
\midrule
swiglu vs.\ bilinear\ (add a gate)          & 0.89 & 0.22 & \textbf{0.17} \\
gelu vs.\ relu2\ (swap elementwise fn)      & 0.97 & 0.89 & \textbf{0.74} \\
\bottomrule
\end{tabular}
\caption{Cross-arm cosine of the reference $\bl$ between matched arms trained
from bit-identical initialisation (median over 12 layers). Adding a gate
relocates $\bl$; swapping the elementwise function does not. Both pairs start
aligned; the divergence is trained. At step 20000 a double-gated arm sits at
$0.72$ from swiglu and $0.18$ from bilinear. Source:
\texttt{ladder\_gate\_reference.py}.}
\label{tab:reloc}
\end{table}

\paragraph{The double-gated arm.}
A third gated arm, with the SiLU on both branches (again bit-identical
initialisation to the SwiGLU), shows where the coupling goes when two branches
carry a threshold rather than one: it splits between them in equal shares. The
two branches sit at alignment ratio $1$, both rest closed, and each carries $\bl$
as its top read principal component at about half the single-gate strength; the
double-gated arm's reference sits at cosine $0.72$ from the SwiGLU and $0.18$ from
the bilinear. Across the three gated arms the pattern is that the reference
coupling distributes over every branch that carries a threshold, in equal shares
by symmetry: one gate takes the coupling and leaves the value branch free, two
gates split it and free neither, and the ungated bilinear product anchors it to
neither branch. This arm is measured at a single training checkpoint.

\paragraph{Which architectures reach which ends.}
Reading the operating point of each arm's gate/read branch at step 150000
(Table~\ref{tab:spectrum}) places the four architectures on the same axis. The
non-gated arms (GELU, squared-ReLU) are almost entirely switches: $88$--$92\%$ of
their units rest off-by-default, but a single read matrix ties key to value, so
they cannot occupy the multiplier end with independent reads. The bilinear arm is
the mirror image: with no off-state, $96\%$ of its units sit at the knee and none
rest closed, so it reaches only the multiplier end. Only SwiGLU carries mass at
both ends ($28\%$ off-by-default, $18\%$ at the knee, the rest in between) with a
wide coupling spread and two independent branches. Reading the median resting
point down the arms gives a monotone march from all-switch (GELU $-2.2$) to
all-multiplier (bilinear $0.0$), with SwiGLU between them. So the gated
architecture is the one that spans the switch--multiplier spectrum with a
decoupled key and value.

\begin{table}[t]
\centering\small
\begin{tabular}{lrrrl}
\toprule
arm & rest (med) & off-by-default & at knee & key/value decoupled \\
\midrule
gelu     & $-2.20$ & 92\% &  1\% & no \\
relu2    & $-1.45$ & 88\% &  3\% & no \\
swiglu   & $-0.84$ & 28\% & 18\% & yes \\
bilinear & $+0.00$ &  0\% & 96\% & yes \\
\bottomrule
\end{tabular}
\caption{Operating-point spectrum across the ladder arms (step 150000, gate/read
branch, pooled over three mid-stack layers). Off-by-default is the fraction of
units with duty $<0.15$; at-knee is duty in $[0.35,0.65]$. Only the gated SwiGLU
carries mass at both ends with independent branches. Source:
\texttt{census\_ladder\_gate\_spectrum.py}.}
\label{tab:spectrum}
\end{table}

\section{Freeing the reference returns a content dimension}
\label{sec:capacity}

The reference coupling is a rank-one overhead shared by every read row: the same
direction $\blh$ enters each row's decomposition, so across the read matrix it
occupies one shared axis rather than per-row capacity. In the non-gated arms that
axis is the read matrix's top principal component (Table~\ref{tab:capacity}):
$\bl$ aligns with the first singular direction of the GELU and squared-ReLU read
matrices at $|\cos| \approx 0.85$, carrying read energy about $17$--$20\times$ the
isotropic share. SwiGLU moves that axis into the gate, where $\bl$ is again the
top component, and leaves the value branch without it: there $\bl$ is only the
matrix's $78$th principal component, at the isotropic floor. The value branch is a
content read matrix with the shared reference removed, and removing it returns the
branch's top principal direction to content. The recovered capacity is real but
small on this measure -- the value branch's effective rank is about $1.5\%$ of $d$
above the GELU read matrix's, matching a GELU matrix with $\bl$ projected out. This
sharpens the anchor's per-row coupling figure, locating the shared slice as the
top read component rather than a diffuse per-row tax; it is a structural
statement, not a performance gain, which this setup does not measure.

\begin{table}[t]
\centering\small
\begin{tabular}{lrrrr}
\toprule
arm : branch & tax & shared/iso & $\bl$ PC\# & PC $|\cos|$ \\
\midrule
gelu : read     & 0.15 & 19.7 & \textbf{1}  & 0.85 \\
relu2 : read    & 0.14 & 16.4 & \textbf{1}  & 0.83 \\
swiglu : gate   & 0.11 & 14.0 & \textbf{1}  & 0.71 \\
swiglu : value  & 0.02 &  1.7 & \textbf{78} & 0.32 \\
bilinear : a/b  & 0.03 & $\sim$2.5 & 2 & 0.40 \\
\bottomrule
\end{tabular}
\caption{Where the shared reference sits in the read matrix (ladder, step
150000, medians over layers). In non-gated arms $\bl$ is the top principal
component of the single read matrix; SwiGLU moves it into the gate and leaves
the value branch's top component free for content. Source:
\texttt{ladder\_content\_capacity.py}.}
\label{tab:capacity}
\end{table}

\section{Discussion}
\label{sec:discussion}

\paragraph{A reason the GLU family is an effective default.}
The spectrum suggests why gating is a safe architectural choice, as a coverage
argument rather than a quality one. A language model uses both thresholded
detectors and multiplicative features; one GLU unit type reaches either, its
position set after the fact by training, whereas the bilinear MLP reaches only
the multiplier end (no off-state) and a single ReLU/GELU unit reaches only the
switch end (a multiplier needs about two units, by the polarisation identity
$ab=\tfrac14[(a{+}b)^2-(a{-}b)^2]$). The gated unit and the bilinear MLP spend
the same per-unit budget -- two reads and a down-projection write, $3d$
parameters -- but only the gated unit covers both ends, because its SiLU sits
on exactly one of the two reads: that read's coupling to $\bl$ can rest the
unit closed, which neither of the bilinear MLP's two linear reads can do, so it
always multiplies. A single ReLU/GELU unit spends less ($2d$: one read, one
write) and is pinned to the switch end for the complementary reason -- it has no
second read to multiply against. This is a capability statement,
not a prediction of lower loss; why our ladder cannot adjudicate one, and where
the loss comparison at this feed-forward geometry is made instead, is in
\S\ref{sec:scope}. The coverage argument speaks to what each architecture can
represent, not to a measured win.

\paragraph{Every trained unit uses both reads.}
One thing the measurements settle is that the switch--multiplier split is not a
split into single-read and two-read units. Trained SwiGLU has independent key and
value reads on essentially every unit ($\cos(\wg,\wv)\approx 0$ throughout),
switch-end units included, so a unit's coupling to the reference sets its
operating point, not whether it uses its second read. That value read earns its
keep across the whole population, not only at the multiplier end, which is the
substance behind the $3d$-versus-$2d$ accounting GLU implementations use to
match a non-gated MLP's width. Oskin's NC-FFN (\S\ref{sec:related}), a GELU
highway with a minority of multiplicative operand units \citep{oskin2026}, ties
GELU, a fully-multiplicative version fails to train, and its operators prefer a
soft, graded regime whose crisp two-operand logic erodes over training (his
\S7.6).
That preference matches the drift toward the soft, near-knee end we observe, with
one difference worth stating: his imposed second operand decays toward
single-operand gating, whereas SwiGLU's value branch stays informative
(\S\ref{sec:scope}), so the soft units here remain genuine two-read multipliers.

\paragraph{The reading extends inside mixture-of-experts experts.}
Each expert in an MoE is itself a gated MLP, and an expert sees only the tokens
routed to it, so its carried reference should be the mean over its \emph{own}
routed tokens rather than a global mean. A pilot on OLMoE-1B-7B (64 experts,
top-8; three mid-stack layers) bears this out with two qualifications. The
per-expert reading holds: within each expert the coupling localises to the gate
branch (gate/value alignment ratio $4.0/2.3/1.7$, weakening with depth as in the
dense soft spot), and each expert's gate rests closed on its \emph{own} routed
reference ($\mathrm{SiLU}(\wg\!\cdot\!\bl^{(e)})<0$). The reference is
expert-conditional but shared-dominated: an expert's routed mean sits at cosine
$0.79$--$0.86$ to the global mean and $0.61$--$0.73$ to other experts' references,
a per-expert tilt on a common reference that grows with depth. Most cleanly, the
routed reference is essentially orthogonal to the expert's router direction
($\cos \approx 0.00$--$0.03$): the computation-stage reference and the
selection-stage routing vector are distinct objects, which makes the separation
from \citet{li2026confadaptive} -- who project the gate weight onto the router
direction as a per-token confidence bias -- an empirical finding rather than only
a conceptual one.

A second pilot on Qwen1.5-MoE-A2.7B (60 routed experts top-4 plus one shared
expert) adds the contrast OLMoE lacks. The shared expert, which processes every
token, reads the most like a dense MLP against the global reference: its
gate/value coupling ratio ($2.5$--$10.6\times$ across the three layers) exceeds
the routed experts' ($1.4$--$3.7\times$), it aligns to the global reference more
than they do, and it rests closed. The routed experts instead couple to their own
tilted references (cosine $0.86$--$0.88$ to the global mean, $0.71$--$0.78$ across
experts), and the routed-reference-versus-router-direction cosine is again small
(here mildly negative, $-0.20$ to $-0.27$, against $\approx 0$ in OLMoE), far from
the value coincidence would give. So across two MoE families the always-on expert
carries the global reference while routed experts specialise to their routed ones,
and the two-level separation holds. These are single-seed pilots on two models;
the depth and cross-family details are not pinned down.

\section{Scope and limitations}
\label{sec:scope}

\begin{itemize}
\item \textbf{Both branches carry content.} What localises to the gate is the
  shared reference slice, and only that. The gate branch carries content as well:
  its row is the coupling $c_n\,\blh$ plus a content remainder $w_\perp$ that is
  about $99\%$ of the row, exactly as in the anchor's per-row anatomy. Freezing
  either branch at its resting value does comparable damage to prediction
  (value-over-gate top-1 flip $0.84$--$1.18$ and $\Delta$NLL $0.5$--$1.8\times$
  across the three layers), as a product of two factors requires.
\item \textbf{The coupling sits in the reference's high-variance subspace.} As
  \S\ref{sec:localize} reports, the localisation survives removing the eight
  massive coordinates (ratio $3$--$13\times$) but narrows toward parity once the
  top sixty-four variance coordinates are gone. The claim is a high-variance-%
  subspace coupling, not a coupling to the reference as a whole.
\item \textbf{The multiplier reading is leading-order.} A knee-resting unit is
  bilinear only to first order in the SiLU expansion; the curvature term and
  finite pre-activation add higher-order structure, so it approximates rather than
  equals a product. The within-model spectrum (Table~\ref{tab:modes}) is a single
  model (Qwen2.5-3B); the cross-architecture span (Table~\ref{tab:spectrum}) is
  the 160M ladder.
\item \textbf{Quality is out of scope, and our ladder cannot speak to it.} The
  ladder is 160M and trained only to a few billion tokens, so it grounds the
  geometric relocation but is far too small and short to adjudicate loss
  differences; we draw no quality conclusion from it, in either direction. The
  loss comparison at this feed-forward geometry is elsewhere \citep{shazeer2020};
  our claim is about where a shared reference sits in the two branches.
\item \textbf{Two seeds on the ladder relocation.} The relocation reproduces at a
  second initialisation: at step 12000 (about $80\%$ of the effect) the
  SwiGLU--bilinear cosine is $0.30$ against the GELU--squared-ReLU control's
  $0.84$, the same dissociation as the first seed ($0.22$ versus $0.89$ at that
  step). The full-horizon replicate and further seeds are not run.
\item \textbf{The reference is input-conditioned and rotates across depth.}
  Inherited from the anchor: $\bl$ shifts with the input distribution (prose to
  code cosine $0.76$--$0.84$, orthogonal for a degenerate repeated token) and
  rotates slowly across the stack. Every statement here is per layer and for the
  reading distribution rather than about a single global direction.
\end{itemize}

\section{Related work}
\label{sec:related}

This note extends \citet{tuomi2026background}: that paper establishes the
coupling $w\cdot\bl$ within a single read matrix, and the result here is that in a
gated MLP the coupling sits in the gate branch while the value branch is free of
it, a two-branch contrast the single-matrix treatment does not make.

Gated linear units were introduced by \citet{dauphin2017}; \citet{shazeer2020}
evaluates GLU variants, including SwiGLU and the bilinear MLP, in the
Transformer feed-forward layer, finding them at least matching GELU and
offering no account of why; the
object here is functional, where each unit sits on the switch--multiplier
spectrum. The multiplier end is the bilinear MLP that \citet{pearce2024} analyse
as a product with no elementwise nonlinearity; our reading places ordinary SwiGLU
units on a continuum whose far end is that bilinear object, indexed by coupling to
$\bl$. The nearest architecture that assigns the two roles to separate unit types
by construction is Oskin's NC-FFN \citep{oskin2026}; we compare its trained
behaviour to the depth-wise drift we measure once the switch--multiplier
coordinate is defined, in \S\ref{sec:discussion}. The global activation shape
gives a different, shared interpolation:
SiLU/Swish move from linear to ReLU through one parameter $\beta$
\citep{ramachandran2017swish}, whereas our coordinate is per unit and set by the
carried reference. The closest mechanism is \citet{li2026confadaptive}, who split
an expert gate weight into a component along the router direction (a per-token
``confidence'' bias that slides a unit along the SiLU curve) and an orthogonal
remainder; we differ in the reference (a carried corpus-mean, not a learned router
vector), the setting (trained dense models, not a proposed MoE architecture), the
coordinate (per unit and static, not per token), and in naming the multiplier
pole. Weight-based taxonomies of gated neurons are adjacent: \citet{gerstner2025}
classify units by input--output weight cosine with a depth trend (a different
cosine, measuring residual read/write alignment) and \citet{gluscope2026} by the
four gate/value sign regimes; neither uses a carried reference. Finally, massive
activations have been read as carried near-constant components with a bias-like
role \citep{sun2024} and MLP units as binary routers \citep{balogh2026}; the
reading here is consistent with both, and \S\ref{sec:localize} identifies the
reference-carrying coordinates as the massive ones.

\section*{Reproducibility}
{\raggedright
Scripts and JSON artifacts (fixed seeds, public checkpoints where applicable):
\texttt{census\_gated\_gate\_structure.py} (Table~\ref{tab:geom}, gate/value
geometry), \texttt{census\_gated\_massive\_control.py} (the massive-coordinate control,
\S\ref{sec:localize}), \texttt{census\_gated\_branch\_xmodel.py} and
\texttt{census\_gated\_uncoupled\_smollm2.py} (Table~\ref{tab:xmodel}),
\texttt{census\_bL\_shape\_xarch.py} (the cross-%
architecture $\bl$-shape measurement, \S\ref{sec:localize}),
\texttt{census\_gated\_coupling\_dist.py} and
\texttt{census\_gated\_uncoupled\_units.py} (Table~\ref{tab:modes}, the
switch--multiplier split and the content-gate characterisation),
\texttt{census\_gate\_bilinear\_fidelity.py} (the behavioural bilinear-fidelity
check, \S\ref{sec:spectrum}),
\texttt{census\_gated\_value\_content.py} (\S\ref{sec:scope}),
\texttt{ladder\_gate\_reference.py} (Table~\ref{tab:reloc}),
\texttt{census\_ladder\_gate\_spectrum.py} (Table~\ref{tab:spectrum}) and
\texttt{ladder\_content\_capacity.py} (Table~\ref{tab:capacity}), with the
feed-forward variants in \texttt{ffn\_variants.py}. The training ladder is the
bit-identical-initialisation four-arm set of \citet{tuomi2026background}'s
companion GLU study; the seed replicate uses \texttt{run\_glu\_matched.py} with a
second initialisation and \texttt{measure\_seed1\_relocation.py}. The
mixture-of-experts pilots are \texttt{census\_moe\_expert\_reference.py} (OLMoE)
and \texttt{census\_moe\_shared\_vs\_routed.py} (Qwen1.5-MoE), \S\ref{sec:discussion}.
Code and data are available at
\url{https://github.com/EvidentSolutions/llm-interp/tree/main/background_glu}
and archived on Zenodo.
\par}

\section*{Use of AI Assistants}
Large language models were used as assistive tools for coding, running
experiments, and drafting text. All research questions, experimental design, and
reported claims were directed and verified by the author.

\bibliographystyle{plainnat}
\bibliography{background_glu}

\end{document}
